% ============================================================
% Collateral Damage in Graph-based Memory Forgetting:
% Quantification, Exploitation, and Defense
% ============================================================
\documentclass[11pt]{article}

% --- Packages ---
\usepackage[utf8]{inputenc}
\usepackage[T1]{fontenc}
\usepackage{microtype}
\usepackage{amsmath,amssymb}
\usepackage{booktabs}
\usepackage{graphicx}
\usepackage{xcolor}
\usepackage{hyperref}
\usepackage{cleveref}
\usepackage{multirow}
\usepackage{caption}
\usepackage{subcaption}
\usepackage[a4paper,margin=2.5cm]{geometry}
\usepackage{natbib}
\usepackage{algorithm}
\usepackage{algorithmic}
\usepackage{tcolorbox}

% --- Custom Commands ---
\newcommand{\attraware}{\textsc{Attr-Aware}}
\newcommand{\bfsprop}{\textsc{BFS}}
\newcommand{\finding}[1]{\vspace{2pt}\noindent\textbf{Finding #1.}~}

% --- Title ---
\title{Collateral Damage in Graph-based Memory Forgetting:\\
Quantification, Exploitation, and Defense}

\author{
  Mingyu Choi \\
  \textit{Doosan Digital Innovation} \\
  \texttt{mingyu1.choi@doosan.com}
}

\date{}

\begin{document}
\maketitle

% ============================================================
\begin{abstract}
Graph-based propagation is a natural approach to invalidating stale facts in AI memory systems: when a fact changes, propagate the update through an entity graph to downstream dependents.
We show that this intuition leads to \textbf{catastrophic collateral damage} that is both quantitatively severe and adversarially exploitable.
On MemoryAgentBench, BFS-based propagation through entity co-occurrence graphs decays 68.9--78.5\% of valid memories across scales (6K--64K facts), with hub entity degrees growing superlinearly ($23 \rightarrow 121 \rightarrow 298$).
More critically, we demonstrate that an attacker can \textbf{weaponize this mechanism}: injecting just 5 fake facts targeting hub entities destroys 57.3\% of valid memories at 32K context.
We further show that this collateral damage \textbf{crosses task boundaries}---BFS propagation triggered during fact consolidation degrades retrieval accuracy on an unrelated task (EM $-5.0$pp, F1 $-18.2$pp).
We propose \textbf{Attribute-aware Selective Propagation}, inspired by directional link failure models in transportation engineering, which defends against both organic and adversarial damage with 78--100\% effectiveness across all experiments.
Graph topology analysis confirms the co-occurrence graph exhibits scale-free properties ($\alpha = 2.30$--$2.48$), providing theoretical grounding from percolation theory for why hub-targeted attacks are structurally devastating.
\end{abstract}

% ============================================================
\section{Introduction}
\label{sec:intro}

AI memory systems that persist and update factual knowledge face a fundamental tension: when stored facts become outdated, the system must invalidate not only the stale fact itself but also dependent knowledge.
Graph-based propagation---traversing entity connections to propagate decay---is the natural solution \citep{mnemosyne_2025, magma_2026}.

\textbf{We discover three critical flaws in this approach.}

First, BFS propagation through entity co-occurrence graphs causes \textbf{catastrophic collateral damage}: at 64K context, a single fact update can cascade through hub nodes to decay 78.5\% of all valid memories.
The root cause is hub entities whose degree grows superlinearly with context size ($23 \rightarrow 298$ from 6K$\rightarrow$64K).

Second, this vulnerability is \textbf{adversarially exploitable}.
An attacker who identifies hub entities in the graph can inject a small number of fake facts (1--5) to trigger the forgetting mechanism, destroying up to 57.3\% of valid memories.
This represents a novel attack vector on AI memory systems that we term \textbf{hub exploitation}.

Third, the damage \textbf{crosses task boundaries}.
BFS propagation triggered during fact consolidation degrades retrieval accuracy on completely unrelated queries, reducing F1 by 18.2 percentage points.

We propose \textbf{\attraware{} Selective Propagation}, a directional propagation method inspired by transportation network link failure models \citep{chen_network_2012}, and demonstrate that it defends against all three threat vectors with 78--100\% effectiveness.

\subsection{Contributions}

\begin{enumerate}
  \item \textbf{Scale trend analysis (C1)}: Collateral damage quantification across 6K$\rightarrow$32K$\rightarrow$64K, revealing superlinear hub degree growth and stable $\sim$78\% \bfsprop{} damage ratio. Graph topology analysis confirms scale-free properties ($\alpha = 2.30$--$2.48$) with heterogeneity parameter $\kappa$ growing superlinearly ($4.6 \rightarrow 32.3$).

  \item \textbf{Cross-task contamination (C2)}: First demonstration that graph-based forgetting in one task degrades retrieval accuracy in an unrelated task (EM $-5.0$pp, F1 $-18.2$pp on Accurate\_Retrieval).

  \item \textbf{Adversarial hub exploitation (C3)}: A novel attack scenario where 5 fake facts destroy 57.3\% of valid memories through weaponized BFS propagation, with scale amplifying vulnerability (5.2\%$\rightarrow$38.1\% damage from a single attack at 6K$\rightarrow$32K). Black-box attacks (general knowledge only) achieve 0\% hit rate, establishing that gray-box or higher knowledge is required.

  \item \textbf{\attraware{} defense (C4)}: Consistent defense effectiveness of 78--100\% across all organic and adversarial damage scenarios. Comparison with degree-capped BFS baselines shows attribute-aware is 80$\times$ more effective than the best structural defense.
\end{enumerate}


% ============================================================
\section{Background and Motivation}
\label{sec:background}

\subsection{AI Memory and Fact Invalidation}

Long-term AI memory systems \citep{mem0_2024, memgpt_2023, amem_2025} store user interactions as structured facts.
When new information contradicts stored facts, the system must (1)~detect the conflict, (2)~invalidate the outdated fact, and (3)~propagate the change to dependent downstream facts.
Step~3 is where graph-based approaches are used, but \textbf{none of the existing works quantitatively analyze the collateral damage of propagation, its adversarial exploitability, or its cross-task effects}.

\subsection{Transportation Network Analogy}

In transportation engineering, \textbf{link failure propagation} studies how road closures cascade through networks \citep{chen_network_2012}.
Hub intersections amplify disruption disproportionately---closing one road at a major interchange can gridlock an entire region.
Critically, transportation research distinguishes between \textit{directional} and \textit{non-directional} propagation: closing Seoul$\rightarrow$Busan affects only Busan-downstream traffic, not Seoul$\rightarrow$Incheon.

Standard BFS propagation in AI memory is non-directional---it spreads invalidation to all connected entities regardless of semantic direction.
Our attribute-aware approach applies the directional principle.

\subsection{MemoryAgentBench}

MemoryAgentBench \citep{memoryagentbench_2026} provides standardized evaluation across four task types: Conflict\_Resolution (FactConsolidation), Accurate\_Retrieval, Test\_Time\_Learning, and Long\_Range\_Understanding.
We use Conflict\_Resolution for scale analysis (C1) and adversarial experiments (C3), and Accurate\_Retrieval for cross-task analysis (C2).


% ============================================================
\section{Collateral Damage at Scale (C1)}
\label{sec:c1}

\subsection{Experimental Setup}

We measure collateral damage by running FactConsolidation memorization under three modes:

\begin{table}[h]
\centering
\small
\begin{tabular}{lll}
\toprule
\textbf{Mode} & \textbf{Conflict Detection} & \textbf{Propagation} \\
\midrule
Baseline & Subject-key match & None \\
\bfsprop{} & Subject-key match & BFS through co-occurrence graph \\
\attraware{} & Subject-key match & Downstream fact chain only \\
\bottomrule
\end{tabular}
\caption{Experimental modes for collateral damage measurement.}
\label{tab:modes}
\end{table}

\textbf{Infrastructure}: Gemma-4-31B-it (LLM), BAAI/bge-m3 1024-dim (Embedding), SQLite + NumPy in-memory, NetworkX graph.

\subsection{Scale Trend Results}

\begin{table}[h]
\centering
\small
\begin{tabular}{lrrrrrl}
\toprule
\textbf{Scale} & \textbf{Valid} & \textbf{BFS Decay} & \textbf{Attr Decay} & \textbf{Reduction} & \textbf{BFS Kill} & \textbf{Max Hub} \\
\midrule
6K   & 344   & 237 (68.9\%) & 51 (14.8\%) & 78.5\% & 143  & 23 \\
32K  & 1,744 & 1,377 (79.0\%) & 275 (15.8\%) & 80.0\% & 1,179 & 121 \\
64K  & 3,414 & 2,680 (78.5\%) & 513 (15.0\%) & 80.9\% & 2,465 & 298 \\
\bottomrule
\end{tabular}
\caption{Collateral damage across scales. BFS Kill = memories with weight $< 0.1$. Max Hub = maximum entity degree.}
\label{tab:scale_trend}
\end{table}

\finding{1} BFS damage ratio plateaus at $\sim$78--79\% across scales, but the \textbf{absolute number of destroyed memories scales linearly}: 143$\rightarrow$1,179$\rightarrow$2,465 complete losses.

\finding{2} Hub entity degree grows superlinearly: $23 \rightarrow 121 \rightarrow 298$ ($\sim$13$\times$ from 6K$\rightarrow$64K), while fact count grows $\sim$7.5$\times$.

\finding{3} \attraware{} damage ratio remains stable at $\sim$15\% regardless of scale.

\subsection{Graph Topology Analysis}

We verify that the entity co-occurrence graph exhibits scale-free properties using maximum likelihood estimation for power-law exponents \citep{clauset_powerlaw_2009}.

\begin{table}[h]
\centering
\small
\begin{tabular}{lrrrrrr}
\toprule
\textbf{Scale} & \textbf{Nodes} & \textbf{$\langle k \rangle$} & \textbf{$k_{\max}$} & \textbf{Gini} & \textbf{$\alpha$ (MLE)} & \textbf{$\kappa$} \\
\midrule
6K  & 414   & 2.51 & 23  & 0.380 & $2.48 \pm 0.09$ & 4.6 \\
32K & 1,732 & 3.00 & 121 & 0.487 & $2.35 \pm 0.04$ & 14.8 \\
64K & 3,187 & 3.26 & 298 & 0.534 & $2.30 \pm 0.03$ & 32.3 \\
\bottomrule
\end{tabular}
\caption{Graph topology across scales. $\alpha$: power-law exponent (MLE, $k_{\min}=2$). $\kappa = \langle k^2 \rangle / \langle k \rangle$: heterogeneity parameter (Molloy--Reed criterion). Gini: degree inequality.}
\label{tab:topology}
\end{table}

\finding{4} All three scales yield $\alpha \in [2.30, 2.48]$, within the classic scale-free range $2 < \alpha < 3$ \citep{barabasi_scalefree_1999}. The heterogeneity parameter $\kappa$ grows superlinearly ($4.6 \rightarrow 32.3$), meaning the network becomes \textit{more} vulnerable to targeted hub attacks at scale---consistent with percolation theory predictions for scale-free networks \citep{cohen_percolation_2000, albert_attack_2000}.

\finding{5} Hub concentration intensifies with scale: the top 1\% of nodes control 7.3\% (6K), 15.5\% (32K), and 19.8\% (64K) of all edges.


% ============================================================
\section{Cross-task Contamination (C2)}
\label{sec:c2}

\subsection{Setup}

We use the Accurate\_Retrieval split (eventqa\_65536) from MemoryAgentBench. Three conditions: (1) Baseline (no propagation), (2) Post-BFS, (3) Post-\attraware{}.

\begin{table}[h]
\centering
\small
\begin{tabular}{lrrrl}
\toprule
\textbf{ARM} & \textbf{EM} & \textbf{F1} & \textbf{Collateral} & \textbf{Affected} \\
\midrule
Baseline     & 15.0\% & 67.1\% & 0  & 0 \\
Post-BFS     & \textbf{10.0\%} & \textbf{48.9\%} & 32 & 3,381 \\
Post-Attr    & 15.0\% & 67.1\% & 1  & 0 \\
\bottomrule
\end{tabular}
\caption{Cross-task contamination. A single BFS propagation from hub ``Debbie'' (degree=715) cascades to 3,381 memories.}
\label{tab:cross_task}
\end{table}

\finding{6} BFS propagation from a single hub entity causes EM $-5.0$pp and F1 $-18.2$pp on \textit{completely unrelated} retrieval queries. \attraware{} preserves full accuracy.


% ============================================================
\section{Adversarial Hub Exploitation (C3)}
\label{sec:c3}

\subsection{Threat Model}

We consider an attacker who can inject text into the memory system's input stream.
The attacker's goal is to maximize destruction of valid memories with minimal input.
We evaluate three attacker knowledge levels:

\begin{itemize}
  \item \textbf{Black-box}: General knowledge only (e.g., ``The capital of France is Berlin'')
  \item \textbf{Gray-box}: Knowledge of which entities are likely hubs (high-degree nodes)
  \item \textbf{White-box}: Access to the fact registry and subject-key generation rules
\end{itemize}

Our primary experiments use white-box targeting (registry-based key lookup) to establish upper-bound damage. We additionally evaluate black-box attacks to establish the lower bound.

\subsection{White-box Results}

\begin{table}[h]
\centering
\small
\begin{tabular}{lrrrr}
\toprule
\textbf{Attack} & \textbf{BFS Damage} & \textbf{BFS Kill} & \textbf{Attr Damage} & \textbf{EM Drop} \\
\midrule
\multicolumn{5}{l}{\textit{6K Context (344 valid memories)}} \\
1 fact & 5.2\% & 0.3\% & 0.0\% & +0.0pp \\
3 facts & 26.7\% & 4.4\% & 0.0\% & $-$10.0pp \\
5 facts & 39.8\% & 17.4\% & 0.3\% & $-$10.0pp \\
\midrule
\multicolumn{5}{l}{\textit{32K Context (1,744 valid memories)}} \\
1 fact & 38.1\% & 1.0\% & 0.0\% & +0.0pp \\
3 facts & 46.0\% & 8.8\% & 1.0\% & +0.0pp \\
5 facts & \textbf{57.3\%} & \textbf{19.3\%} & 1.0\% & +0.0pp \\
\bottomrule
\end{tabular}
\caption{White-box adversarial attack results.}
\label{tab:adversarial}
\end{table}

\finding{7} A single fake fact destroys 5.2\% of valid memories at 6K and 38.1\% at 32K---a $7.3\times$ amplification from scale alone.

\finding{8} Five fake facts at 32K destroy 57.3\% of valid memories and cause 19.3\% complete loss (weight $< 0.1$). The system's own forgetting mechanism becomes a weapon.

\subsection{Black-box Attack Results}

\begin{table}[h]
\centering
\small
\begin{tabular}{lrrr}
\toprule
\textbf{Attack Type} & \textbf{Facts} & \textbf{Hit Rate} & \textbf{BFS Damage} \\
\midrule
White-box (registry) & 5 & 100\% & 39.8\% \\
Black-box (general knowledge) & 5 & 0\% & 0.0\% \\
\bottomrule
\end{tabular}
\caption{Black-box vs white-box attack comparison (6K). Black-box attacks fail to trigger conflict detection because general-knowledge facts do not match the subject-key regex pattern.}
\label{tab:blackbox}
\end{table}

\finding{9} Black-box attacks achieve 0\% hit rate, demonstrating that the subject-key matching mechanism provides implicit defense against naive attackers. Gray-box or higher knowledge of the system's fact structure is required for successful exploitation.

\subsection{Propagation Depth Analysis}

\begin{table}[h]
\centering
\small
\begin{tabular}{lrrrrr}
\toprule
\textbf{Depth} & \textbf{Damage} & \textbf{Kill} & \textbf{Hop-1} & \textbf{Hop-2} & \textbf{Hop-3+} \\
\midrule
1 & 18.6\% & 0.3\% & 232 & --- & --- \\
2 & 39.8\% & 17.4\% & 232 & 490 & --- \\
3 & 48.3\% & 32.9\% & 232 & 490 & 605 \\
4 & 55.8\% & 44.8\% & 232 & 490 & 1,131 \\
\bottomrule
\end{tabular}
\caption{Propagation depth vs damage (6K, 5-fact attack). Memory counts indicate the number of memories affected at each hop distance.}
\label{tab:depth}
\end{table}

\finding{10} Hop-1 damage is constant (232 memories), but distant-hop impact grows rapidly (Hop-3: 605, Hop-4: 1,131). Kill rate \textit{accelerates} with depth due to cumulative decay ($0.3\% \rightarrow 44.8\%$).


% ============================================================
\section{Attribute-aware Selective Propagation (C4)}
\label{sec:c4}

\subsection{Method}

When fact $f$ with subject\_key \texttt{Entity:Attribute} is invalidated:
\begin{enumerate}
  \item Extract \textit{object entities} from $f$ (non-subject entities)
  \item Find downstream facts where object entities are subjects
  \item Apply decay: $w_{\text{new}} = w_{\text{curr}} \times (\delta^h)$ with in-degree adjustment
  \item Recurse for depth $> 1$
\end{enumerate}

This restricts propagation to genuine causal chains rather than all graph neighbors.

\subsection{Defense Comparison}

\begin{table}[h]
\centering
\small
\begin{tabular}{lrrrr}
\toprule
\textbf{Defense} & \textbf{Damage} & \textbf{Kill} & \textbf{Post-EM} & \textbf{Post-F1} \\
\midrule
\bfsprop{} (no cap) & 39.8\% & 17.4\% & 10.0\% & 13.3\% \\
Degree-cap=50 & 39.8\% & 17.4\% & 10.0\% & 13.3\% \\
Degree-cap=20 & 33.7\% & 17.4\% & 10.0\% & 13.3\% \\
Degree-cap=10 & 28.2\% & 10.8\% & 10.0\% & 13.3\% \\
Degree-cap=5 & 24.1\% & 8.1\% & 10.0\% & 13.3\% \\
\textbf{\attraware{}} & \textbf{0.3\%} & \textbf{0.0\%} & \textbf{20.0\%} & \textbf{23.3\%} \\
\bottomrule
\end{tabular}
\caption{Defense mechanism comparison (6K, 5-fact attack). Degree capping reduces damage but is insufficient; \attraware{} is $80\times$ more effective than the best structural defense.}
\label{tab:defense}
\end{table}

\finding{11} Degree capping reduces damage from 39.8\% to 24.1\% (cap=5), but remains $80\times$ worse than \attraware{} (0.3\%). Structural defenses cannot substitute for semantic propagation.

\subsection{Effectiveness Across All Experiments}

\begin{table}[h]
\centering
\small
\begin{tabular}{lrrr}
\toprule
\textbf{Experiment} & \textbf{BFS Damage} & \textbf{Attr Damage} & \textbf{Defense Rate} \\
\midrule
CD 6K & 68.9\% & 14.8\% & 78.5\% \\
CD 32K & 79.0\% & 15.8\% & 80.0\% \\
CD 64K & 78.5\% & 15.0\% & 80.9\% \\
Cross-task (F1) & $-$18.2pp & +0.0pp & 100.0\% \\
Adv 6K / 1f & 5.2\% & 0.0\% & 100.0\% \\
Adv 6K / 5f & 39.8\% & 0.3\% & 99.3\% \\
Adv 32K / 1f & 38.1\% & 0.0\% & 100.0\% \\
Adv 32K / 5f & 57.3\% & 1.0\% & 98.2\% \\
\bottomrule
\end{tabular}
\caption{Defense effectiveness across all experiments.}
\label{tab:defense_all}
\end{table}


% ============================================================
\section{Related Work}
\label{sec:related}

\textbf{Memory systems.}
Mem0 \citep{mem0_2024}, MemGPT/Letta \citep{memgpt_2023}, A-Mem \citep{amem_2025}, and MAGMA \citep{magma_2026} implement memory management but do not analyze propagation damage or adversarial vulnerability.

\textbf{Temporal decay.}
Mnemosyne \citep{mnemosyne_2025} uses Ebbinghaus-inspired decay \citep{ebbinghaus_1885}. Our work is complementary---we address \textit{which} memories to decay, not \textit{how fast}.

\textbf{Adversarial attacks.}
Prompt injection \citep{perez_prompt_injection_2022} and data poisoning \citep{data_poisoning_2021} attack language models directly. Our hub exploitation is a novel vector that weaponizes the system's own memory management mechanism.

\textbf{Network resilience.}
Scale-free networks are robust to random failures but fragile to targeted hub attacks \citep{albert_attack_2000, barabasi_scalefree_1999}. Percolation theory \citep{cohen_percolation_2000, molloy_reed_1995} and transportation network resilience \citep{chen_network_2012} provide the theoretical foundation for our analysis. We are the first to apply these principles to AI memory invalidation.


% ============================================================
\section{Discussion}
\label{sec:discussion}

\subsection{Implications for AI Memory Safety}

Our results establish \textbf{memory graph exploitation} as a distinct attack surface in AI systems, separate from prompt injection and data poisoning. Key implications:
\begin{enumerate}
  \item Any system using graph-based forgetting is vulnerable to hub exploitation
  \item The attack is low-cost: 1--5 injected facts can destroy 38--57\% of stored knowledge
  \item The vulnerability worsens with scale: more stored knowledge means larger blast radius
  \item \attraware{} propagation provides a practical defense without significant overhead
\end{enumerate}

\subsection{Why Hub Entities Are Fundamental}

Hub entities reflect genuine reality---the United States really does connect to thousands of facts. The problem is that \textbf{co-occurrence $\neq$ causal dependency}. \attraware{} propagation resolves this by distinguishing genuine dependency chains from coincidental co-occurrence.

\subsection{Limitations}

\begin{enumerate}
  \item \textbf{Entity extraction}: Rule-based extraction treats ``Vladimir Putin'' and ``Putin'' as separate entities. Stronger entity resolution would improve both modes.
  \item \textbf{Single benchmark}: Results from MemoryAgentBench only. Replication on other benchmarks would strengthen findings.
  \item \textbf{Statistical variance}: Multi-run statistics (7 runs at 6K) show large coefficient of variation ($\sim$150\%) across contexts, with BFS vs \attraware{} difference not reaching significance ($p = 0.17$, Cohen's $d = 0.92$). This reflects genuine variability in hub entity structure across contexts rather than method instability.
  \item \textbf{White-box assumption}: Primary adversarial results assume white-box access. Black-box attacks failed entirely (0\% hit rate), suggesting the threat is bounded by attacker knowledge.
  \item \textbf{32K EM floor}: Baseline EM at 32K is 0--4\%, limiting observable EM changes from adversarial attacks at that scale.
\end{enumerate}


% ============================================================
\section{Conclusion}
\label{sec:conclusion}

We present a comprehensive analysis of collateral damage in graph-based memory forgetting across three dimensions: scale, cross-task contamination, and adversarial exploitation.
BFS propagation through entity co-occurrence graphs is \textbf{catastrophically destructive} (78\% damage), \textbf{cross-task toxic} (F1 $-18.2$pp on unrelated retrieval), and \textbf{adversarially weaponizable} (5 fake facts $\rightarrow$ 57.3\% memory destruction).

\attraware{} Selective Propagation, inspired by directional link failure models in transportation engineering, provides consistent defense across all threat vectors (78--100\% effectiveness) with negligible overhead.
Graph topology analysis confirms the underlying mechanism: co-occurrence graphs are scale-free ($\alpha \approx 2.3$--$2.5$), making them structurally robust to random failures but catastrophically vulnerable to targeted hub attacks---a prediction from percolation theory that our empirical results validate.

The key insight: \textbf{propagation must follow causal dependency chains, not structural co-occurrence}---a principle well-established in transportation network resilience but previously unapplied to AI memory systems.

% ============================================================
% Appendices
% ============================================================
\appendix

\section{Experimental Configuration}
\label{app:config}

\begin{table}[h]
\centering
\small
\begin{tabular}{ll}
\toprule
\textbf{Component} & \textbf{Specification} \\
\midrule
LLM & google/gemma-4-31B-it (max\_tokens=10, temp=0.0) \\
Embedding & BAAI/bge-m3 (dim=1024) \\
Vector search & NumPy cosine similarity, RRF fusion ($k=60$) \\
Graph & NetworkX entity co-occurrence (DiGraph) \\
Storage & SQLite in-memory \\
Benchmark & ai-hyz/MemoryAgentBench (ICLR 2026) \\
Context sizes & 6K, 32K, 64K (FactConsolidation), 65K (EventQA) \\
Propagation & depth=2, decay\_per\_hop=0.5 \\
Adversarial & 1/3/5 fake facts $\times$ BFS/Attr $\times$ 6K/32K \\
\bottomrule
\end{tabular}
\caption{Experimental configuration.}
\label{tab:config}
\end{table}


\section{Hub Concentration Analysis}
\label{app:hub}

\begin{table}[h]
\centering
\small
\begin{tabular}{lrrrr}
\toprule
\textbf{Scale} & \textbf{Top 1\%} & \textbf{Top 5\%} & \textbf{Top 10\%} & \textbf{Top 20\%} \\
\midrule
6K  & 7.3\%  & 18.2\% & 27.3\% & 42.4\% \\
32K & 15.5\% & 29.8\% & 39.2\% & 53.0\% \\
64K & 19.8\% & 34.6\% & 43.8\% & 56.8\% \\
\bottomrule
\end{tabular}
\caption{Hub entity concentration: percentage of total edges controlled by top-$N$\% nodes. Inequality increases with scale.}
\label{tab:hub_concentration}
\end{table}


\section{Multi-run Statistical Analysis}
\label{app:stats}

We run the 6K adversarial attack experiment 7 times to assess reproducibility.

\begin{table}[h]
\centering
\small
\begin{tabular}{llrrl}
\toprule
\textbf{Attack} & \textbf{Mode} & \textbf{Damage (mean$\pm$std)} & \textbf{Cohen's $d$} & \textbf{$p$-value} \\
\midrule
1 fact & \bfsprop{} & $1.6 \pm 2.5$ & \multirow{2}{*}{0.95} & \multirow{2}{*}{0.126} \\
       & \attraware{} & $0.0 \pm 0.1$ & & \\
\midrule
3 facts & \bfsprop{} & $8.9 \pm 13.8$ & \multirow{2}{*}{0.93} & \multirow{2}{*}{0.169} \\
        & \attraware{} & $-0.1 \pm 0.4$ & & \\
\midrule
5 facts & \bfsprop{} & $13.3 \pm 20.6$ & \multirow{2}{*}{0.92} & \multirow{2}{*}{0.171} \\
        & \attraware{} & $-0.1 \pm 0.7$ & & \\
\bottomrule
\end{tabular}
\caption{Multi-run statistics (6K, 7 runs). Large effect sizes (Cohen's $d > 0.8$) but high variance yields $p > 0.05$. The CV $\sim$150\% reflects genuine variability in hub entity structure across contexts.}
\label{tab:multirun}
\end{table}


% ============================================================
\bibliographystyle{plainnat}
\bibliography{references}

\end{document}
